% template.tex -- the PRISMA-ScR skeleton that sdr.manuscript renders into main.tex.
% Markers: @@TABLE:<name>@@ -> \input of a generated table; @@FIGURE:<name>@@ -> the figure if
% figures/<name>.pdf exists, else a red placeholder; @@NUMBERS@@ -> the macro file; @@RUN_DATE@@,
% @@LAST_SEARCH_DATE@@, @@PROTOCOL_VERSION@@, @@GIT_COMMIT@@, @@REGISTRATION_DOI@@, @@PENDING_LIST@@.
% Rule (PROTOCOL.md section 13): the Abstract and Results carry no bare digit; every number is a
% \num<Name> or \numword<Name> macro from tables/numbers.tex. Never edit main.tex; edit this file.
% Target venue: Systematic Reviews (BMC) -- it publishes protocol-driven scoping reviews with
% code and has no page limit for the appendices this review needs; BMJ Open and JBI Evidence
% Synthesis are the named alternatives and the headings below are theirs too.
\documentclass[11pt,a4paper]{article}
\usepackage[utf8]{inputenc}
\usepackage[T1]{fontenc}
\usepackage{lmodern}
\usepackage{microtype}
\usepackage[margin=25mm]{geometry}
\usepackage{amsmath}
\usepackage{array}
\usepackage{longtable}
\usepackage{booktabs}
\usepackage{tabularx}
\usepackage{graphicx}
\usepackage{xcolor}
\usepackage{caption}
\usepackage{float}
\usepackage{setspace}
\usepackage{lineno}
\usepackage{fancyvrb}
\usepackage[numbers,sort&compress]{natbib}
\usepackage[hidelinks,breaklinks]{hyperref}
\urlstyle{tt}
\setlength{\tabcolsep}{3pt}
\captionsetup{font=small,labelfont=bf}
\linenumbers

% ---- macros the generated files rely on
\newcommand{\pending}[1]{\textcolor{red}{[#1: pending]}}
\newcommand{\pendingtable}[2]{\par\noindent\textcolor{red}{[Table \texttt{\detokenize{#1}} pending: \detokenize{#2}]}\par}
\newcommand{\pendingfigure}[1]{\par\noindent\fbox{\parbox{0.95\linewidth}{\centering\textcolor{red}{[Figure \texttt{\detokenize{#1}} pending: figures/\detokenize{#1}.pdf not generated]}}}\par}
\newcommand{\tabref}[1]{Table~\ref{tab:#1}}
\newcommand{\figref}[1]{Figure~\ref{fig:#1}}
\newcommand{\yearRange}{2000--2026}
\newcommand{\runDate}{@@RUN_DATE@@}
\newcommand{\lastSearchDate}{@@LAST_SEARCH_DATE@@}
@@NUMBERS@@

\title{Publicly described datasets of human spine computed tomography and projection radiography, \yearRange: a scoping review with programmatic access verification and a download recipe}
\author{Gregory Schwing$^{1}$ and the review team\\[2pt]
\small $^{1}$Department of Neurosurgery, Wayne State University School of Medicine, Detroit, MI, USA}
\date{Protocol version @@PROTOCOL_VERSION@@; search run \runDate; rendered from commit \texttt{@@GIT_COMMIT@@}}

\begin{document}
\maketitle

\begin{abstract}
\noindent\textbf{Background.} Level identification from a limited field of view needs images with correct level names, a known annotation origin and labelled variant anatomy (a sixth lumbar vertebra, a thirteenth rib-bearing vertebra, a lumbosacral transitional vertebra). No enumeration of spine CT and radiograph datasets states, for each, how it was annotated, whether variants are labelled and whether it downloads today.
\textbf{Objectives.} To map publicly described human spine CT and projection radiography datasets (\yearRange), annotated or not; to chart annotation type and origin, variant labelling, licence and access route; to score each on a predefined rubric; to produce a download recipe.
\textbf{Eligibility criteria.} Human spine image datasets in CT or projection radiography, at least ten subjects, described in an article, preprint, challenge page or repository record; MRI-only, non-human, phantom-only, synthetic and feature-only releases excluded.
\textbf{Sources of evidence.} \numwordSourcesPlanned{} literature indexes, data repositories and figure-mining routes queried by code with verbatim strings recorded; known items and one hop of citation chasing.
\textbf{Charting methods.} A frozen JSON-schema form with a verbatim quote per charted value; a ten-item rubric scored by rule; an HTTP probe of every access URL on two dates.
\textbf{Results.} The search identified \numIdentified{} records; \numAfterMerge{} remained after de-duplication, \numScreened{} were screened and \numwordDatasets{} datasets were included (\numwordDatasetsCT{} CT, \numwordDatasetsXR{} radiograph, \numwordDatasetsFluoro{} fluoroscopy). \numwordLevelNamesYes{} released per-vertebra level names, \numwordCountingInferred{} counting-inferred names; \numwordVariantExplicit{} labelled a variant per subject. \numwordOpenDownload{} were openly downloadable, \numwordRegistration{} needed registration, \numwordDUA{} a data-use agreement, \numwordRequestToAuthors{} a request to the authors and \numwordUnavailable{} were unavailable; of \numwordDescribedPublic{} described as public, \numPctDescribedPublicDownloads{} downloaded. Median quality total was \numQualityMedian{} (IQR \numQualityQOne{} to \numQualityQThree); \numwordFitLevelId{} datasets were fit for level identification and \numwordFitStitching{} for whole-spine or thoracic work.
\textbf{Conclusions.} The map, scores, probe results and a recipe for every downloadable dataset are released as code and data, so the enumeration can be rerun on any date.
\end{abstract}

\noindent\textbf{Keywords:} spine; computed tomography; radiography; fluoroscopy; datasets; scoping review; annotation; lumbosacral transitional vertebra; data availability.

\section{Introduction}

\subsection{Rationale}\label{sec:rationale}
Vertebral level identification from a limited field of view --- the problem this group works on in fluoroscopy and radiography --- depends on training and test data in which the level names are right, the origin of each label is known, and variant anatomy (a sixth lumbar vertebra L6, a thirteenth rib-bearing vertebra T13, a lumbosacral transitional vertebra LSTV) is labelled rather than silently normalised. The group already uses CTSpinoPelvic1K~\citep{schwing2026ctspinopelvic1k}, VerSe~\citep{sekuboyina2021verse}, TotalSegmentator~\citep{wasserthal2023totalsegmentator}, the NHANES~II films~\citep{long2005nhanes}, BUU-LSPINE~\citep{klinwichit2023buu}, VinDr-SpineXR~\citep{nguyen2021vindr}, AASCE~\citep{wu2017boostnet} and a literature-mined fluoroscopy collection, and has repeatedly found a relevant dataset only after a model had been trained without it.

No published enumeration of spine CT and radiograph datasets states, for each one, how its annotations were produced, whether variant anatomy is labelled, what the licence permits, and whether the data can be downloaded today. Existing lists (dataset roadmaps, challenge overviews, curated lists, Papers with Code) are unsystematic, undated and do not verify access. Whole-spine radiographs and thoracic-spine images are the scarcest class and the most valuable for frame-stitching work, so a review that finds even a few more matters. A scoping review is the right form because the question is ``what exists and what is it like'', not ``what is the effect of X''~\citep{peters2020scoping,peters2020jbi}. The availability question is empirical and changes over time, so it is answered by an HTTP probe on a recorded date, not by reading the paper.

\subsection{Objectives}\label{sec:objectives}
Population, concept and context: images of the human spine (any region, any age); a dataset --- a named or nameable, bounded collection of subject-level images, with or without annotations, that has a description someone else can read and is intended for reuse; CT (including cone-beam CT) and projection radiography (plain radiograph, fluoroscopy or C-arm, DXA vertebral fracture assessment, EOS biplanar), described in a record dated \yearRange, any language, any country.

Primary questions, each answered per dataset and in aggregate: (RQ1) what publicly described datasets exist, with their size, region, projection and format; (RQ2) how each is annotated and where the annotation came from (expert manual, expert-corrected model, model generated, counting inferred, report derived, other); (RQ3) whether L6, T13 and LSTV are labelled, with counts where reported; (RQ4) under what licence each is released; (RQ5) what the access route is today, established by a recorded probe. Secondary objectives: a predefined ten-item quality score and a fixed ``fit for level identification'' rule (SO1); pooled counts and, where the plan permits, illustrative pooled prevalence and agreement (SO2); a download recipe, a manifest and a split stub (SO3); coverage of the authors' compute grid (SO4). The review is first a review with predefined methods; the recipe is a product of the charting records.

\section{Methods}

\subsection{Protocol and registration}\label{sec:protocol}
The protocol (version @@PROTOCOL_VERSION@@) with its companion files --- the information-source specification with every verbatim query string, the charting schema, the quality rubric, the statistical analysis plan and the PRISMA-ScR checklist --- was frozen after a count-only pilot search and before the full search, tagged in the repository and registered on the Open Science Framework (@@REGISTRATION_DOI@@). The review is reported under PRISMA-ScR~\citep{tricco2018prisma} with the structure of the JBI chapter on scoping reviews~\citep{peters2020jbi} and Cochrane practice for search reporting and pre-specified analysis~\citep{higgins2023cochrane}. Every deviation is logged with its date, reason and effect (Appendix~\ref{app:deviations}). Registration was chosen over PROSPERO, which does not accept scoping reviews.

\subsection{Eligibility criteria}\label{sec:eligibility}
The criteria, their decision rules and the rationale for each are fixed in the protocol and reproduced in \tabref{eligibility}. In brief: a dataset of human spine images in CT, CBCT, plain radiography, fluoroscopy, DXA vertebral fracture assessment or EOS, with at least ten subjects (or, when the subject count is unknown, at least ten images or volumes), described in a journal article, preprint, challenge page or data-repository record dated \yearRange, in any language, with or without annotations. Collections assembled from published figures are eligible and flagged figure-derived; datasets derived from another public dataset are eligible and flagged derived with the parent recorded. Excluded are MRI-only, ultrasound-only and nuclear-medicine-only releases; non-human subjects; phantom, cadaver or specimen-only images; purely synthetic images with no public parent; a stated subject count below ten; releases without subject-level images; and records that are not a dataset (a method paper using a private dataset, a review, a software release without images). A method paper that merely uses an already public dataset and releases nothing new is coded use-only and counted against that dataset, not as an included record. Several records may describe one dataset; the unit of analysis is the dataset and the unit of screening is the record.

@@TABLE:eligibility@@

\subsection{Information sources}\label{sec:sources}
\numwordSourcesPlanned{} search sources were specified: seven literature indexes (PubMed, Europe PMC, OpenAlex, Semantic Scholar, arXiv, bioRxiv/medRxiv, Crossref), fourteen data repositories and indexes (Zenodo, Figshare, Dryad, OSF, Mendeley Data, DataCite, PhysioNet, TCIA/NBIA, Kaggle, Hugging Face, grand-challenge.org, GitHub, Papers with Code, OpenDataLab), eight figure-mining routes that the group's fluoroscopy corpus was built from (NLM Open-i, MultiCaRe~\citep{multicare2023}, Open-PMC-18M~\citep{openpmc18m}, Wikimedia Commons, J-STAGE, KoreaMed Synapse, SciELO, Redalyc) and one supplementary venue sweep (Springer Nature Meta API over \emph{Scientific Data}). Unpaywall was used only to locate full text at charting. Each source's endpoint, authentication, rate ceiling, paging cap and index coverage are tabulated in the protocol and, with the query date and hit counts, in \tabref{sources}. Sources considered and excluded, with the reason: Google Dataset Search and Google Scholar (no API; terms of service forbid automated querying); YouTube, Instagram, Reddit, Flickr and Facebook (media platforms with no subject-level medical-image datasets carrying a stable record; a dataset announced there is indexed by at least one source above); Openverse (aggregates Commons, which is searched directly); Common Crawl (raw web text, not a dataset index); patents (drawings and claims, not released image collections). Other methods: a frozen known-item set, one hop of backward and forward citation chasing through OpenAlex for every included source record, and the dataset tables of review articles retained for chasing.

@@TABLE:sources@@

\subsection{Search}\label{sec:search}
Every query is the conjunction of three concept blocks --- spine terms, modality terms and dataset terms --- translated into each source's syntax; sources without boolean syntax receive a fixed grid of thirty-five strings or a fixed single-term list, and sources with no query interface are harvested in full and filtered locally with the frozen prescreen regular expressions, which is recorded as the query. The blocks were assembled from the titles of the known-item set and the group's earlier ad-hoc searches; the count-only pilot search decided one thing (Europe PMC strings are fielded) and harvested no terms. Year limits are in the string where the source supports a date field and applied locally otherwise; there is no language limit. For every query execution the source, query identifier, verbatim string as sent, endpoint, parameters with secrets removed, retrieval timestamps, reported total, number retrieved, pages, paging cap, date-filter mode and credential mode (by name, never value) are recorded in \texttt{queries.json}, from which Appendix~\ref{app:search} is generated; every HTTP attempt is logged and every response body cached by hash so the run can be replayed offline. The full search ran once on \runDate{} and the run directory is frozen; a search update before submission is a new run directory whose records are de-duplicated against the frozen run and screened, charted and probed by the same procedure. Known-item recall (the share of the frozen set the search retrieved) is reported with a Wilson interval as a check on the search, not on the review, because the known items enter through the other-methods arm regardless.

\subsection{Selection of sources of evidence}\label{sec:selection}
Record identity is one canonical string (DOI, else PMID, PMCID, arXiv identifier, canonical repository URL, or the source's native identifier) hashed once; cross-source merging proceeds by six recorded rules (exact DOI; PMID or PMCID; arXiv--DOI link; Zenodo concept DOI; normalised title with first-author surname and year within one; repository--article relations), with the rule that fired counted in the flow. A rule-based prescreen applies the frozen regular expressions to title, abstract and description and records pass, fail (naming the block) or uncertain; repository records are not required to match the dataset block, and records from the non-English routes are never prescreened. A stratified audit sample of failing records (by failing block and source class) is screened blind to the prescreen outcome and the inclusion rate per stratum is reported with the implied number missed. Every passing or uncertain record is screened on title, abstract or description by one reviewer against the coded criteria, and the decision, primary code, all codes and a reason are recorded per record. A second reviewer independently screens a stratified random sample of at least one hundred records or ten percent, and every record that reaches description review; agreement is reported as Cohen's $\kappa$~\citep{cohen1960kappa} on the binary decision with prevalence-adjusted alternatives~\citep{byrt1993pabak,gwet2008ac1}, and disagreements are resolved by discussion. Single-reviewer screening of the full set is a stated departure from JBI, made because of the expected five-figure record count. An LLM-assisted step is used and named as such: a fixed model at temperature zero proposes a decision and codes for each record before the human decides; every prompt and output is stored, the human decision is recorded independently, and the agreement between the suggestion and the decision is reported separately and never combined with the human $\kappa$. At description review both reviewers read the dataset's own documentation for every included or unclear record; full text is obtained through open-access links or Unpaywall, and records whose full text could not be obtained are charted from the abstract and repository record and counted as such.

\subsection{Data charting process}\label{sec:charting}
One JSON record per dataset, validated against the frozen schema on every build. Five datasets spanning modalities and routes were charted by both reviewers as a pilot before the main charting. One reviewer charts every dataset; as soon as the dataset list is fixed a random sample of at least twenty-five datasets or twenty percent is re-charted blind by the second reviewer, before that reviewer verifies any dataset, and only then does the second reviewer verify every field of every dataset against its quotes. Every charted value not computed by code carries a verbatim sentence with its source and locator; licence, file list, sizes, checksums, DOI, version and access fields are pulled by code from the repository APIs, and the record says which value came from which API. Datasets maintained by a review author are charted, verified and scored by an author who did not create them. The review sends no access requests; a single e-mail to a corresponding author is permitted only to resolve a contradictory description and is recorded.

\subsection{Data items}\label{sec:items}
The fields of the charting form, their definitions and allowed values are generated from the schema into \tabref{items}. The operational definitions that matter most: a \emph{subject} is one person; \emph{described as public} is a reviewer-charted value requiring a quote that the images (not code or supplements) are released, charted before the probe and never updated from it, with an independent regex flag reported beside it; \emph{annotation origin} is one of expert-manual, expert-corrected-model, model-generated, counting-inferred (level names assigned by counting from a landmark rather than by reading each vertebra), report-derived, crowd-or-trainee, caption-derived, unknown or none; a variant is \emph{labelled} when it is present per subject in the released labels or metadata or counted in the text, with the mechanism, the ascertainment, the definition class and whether the sample was enriched for it recorded; \emph{per-vertebra level names} is yes, partial or no.

@@TABLE:items@@

\subsection{Availability probe}\label{sec:probe}
Each dataset's access URL --- the URL the description gives for the images, or one resolved from a named repository or a data DOI, with the basis recorded --- is probed by code on two dates at least forty-eight hours apart, routing repository hosts to their APIs and everything else to a HEAD request followed by a ranged GET of at most sixty-four kibibytes; no dataset file is ever downloaded. Twenty-five ordered rules classify each probe as open-download, registration, data-use-agreement, request-to-authors, unavailable, commercial, auth-required or unresolved; the final route is the classification the two probes agree on, with a third probe on disagreement and a verifier's recorded decision where the machine could not decide. The share of datasets described as public that actually download is defined once: the denominator is the reviewer-charted value, the primary numerator is a final route of open-download, and a second row reports file URLs that answered on both probes.

\subsection{Critical appraisal of individual sources of evidence}\label{sec:appraisal}
Scoping reviews do not require appraisal; this one scores each dataset because a usable ranking is a stated objective. The instrument is a ten-item rubric (documentation; annotation origin; annotation validation; variant-anatomy coverage; metadata completeness; licence clarity; access verified by probe; versioning and identifier; split definition and patient-level de-duplication; image-quality reporting), each item scored zero, one or two by a rule over charted fields only, so that every score traces to a quote or an API value; unannotated datasets score zero on the two annotation items by rule and are reported in their own stratum. Inter-rater reliability comes from the blind re-charting sample: per-item weighted $\kappa$ and ICC(2,1) on the total~\citep{shrout1979icc}. A boolean ``fit for level identification'' is computed by four fixed rules (level names released; their origin expert and not counting-inferred; a route among open-download, registration and data-use-agreement; at least twenty subjects, a threshold that is arbitrary and stated as such), with derived booleans for training use, variant coverage and whole-spine or thoracic coverage. Scores never influence inclusion.

\subsection{Synthesis of results}\label{sec:synthesis}
The statistical analysis plan fixes every table, figure, estimator and eligibility rule; every number in this manuscript is generated by code into a macro file and none is typed. Flow counts per source and in total are asserted by an identity the build checks. Descriptive counts and totals of subjects and images are reported by modality, region, projection, annotation type and origin, access route, licence tier and year, with at most one dataset per parent family in any total and figure-derived and estimated counts in their own rows, labelled ``reported subjects summed across datasets (persons may recur)''. Availability is reported as the route distribution, the described-as-public share, probe agreement between dates, and route by year band with an exploratory trend test. The quality total is described and related to modality and access route (Kruskal--Wallis with $\varepsilon^{2}$, Dunn--Holm post hoc), reported three times --- the full total, without the probe item, and without the three repository-tied items --- and to year (Spearman $\rho$ with a Fisher-$z$ interval~\citep{bonett2000spearman}).

The JBI scoping design maps evidence and does not pool; this review keeps that frame and reports variant prevalence and inter-rater agreement as per-dataset tables with exact intervals~\citep{clopper1934} as the primary output. A pooled summary is produced only as a clearly labelled secondary, illustrative summary, only when at least five eligible datasets exist per variant (subject-level counts, unselected sampling, stated ascertainment and definition, verified, one per parent family), by a logit random-intercept model fitted by adaptive Gauss--Hermite quadrature, with DerSimonian--Laird~\citep{dersimonian1986} and REML with the Hartung--Knapp--Sidik--Jonkman interval~\citep{inthout2014hksj} as sensitivity analyses, $I^{2}$ with a Q-profile interval~\citep{higgins2002} as a description, and the prediction interval as the headline; the Freeman--Tukey transform is not used because its back-transform misleads when sample sizes span orders of magnitude~\citep{schwarzer2019}. Datasets enriched for variants form their own stratum and never enter the pooled estimate. Agreement is pooled only for $\kappa$ groups with equal category counts and reported intervals; Dice, ICC and millimetre errors are summarised as medians. Risk of bias in studies, certainty of evidence and reporting bias are not applicable to a review of datasets and are not assessed. Tests are exploratory, two-sided at five percent, with Wilson intervals~\citep{wilson1927} on proportions. Software: Python with numpy, scipy, pandas, statsmodels and matplotlib~\citep{harris2020numpy,virtanen2020scipy,seabold2010statsmodels}, versions recorded with the git commit.

\section{Results}

\subsection{Selection of sources of evidence}\label{sec:flow}
Of the \numwordSourcesPlanned{} planned sources, \numwordSourcesRun{} returned records through \numQueries{} queries whose requests went out; \numwordQueriesNotSent{} planned queries were not sent, \numwordSourcesNotSearched{} of the planned sources therefore not being searched at all in this version (\tabref{sources} names them with the recorded reason; verbatim strings in Appendix~\ref{app:search}). The search identified \numIdentified{} records; \numAfterMerge{} remained after cross-source de-duplication, \numScreened{} were screened on title, abstract or description, \numSoughtFullText{} were sought for description review, and \numwordIncludedRecordsSearch{} records were included through the search arm; \numwordKnownItems{} known items and the citation-chasing candidates yielded \numwordIncludedRecordsOther{} further included records through the other-methods arm (\figref{flow}). Known-item recall was \numRecallKnownItems{} (\numwordKnownRetrieved{} of \numwordKnownRecallDenominator). On the dual-screened sample of \numDualSample{} records, binary agreement was $\kappa$ = \numKappaDual{} (\numKappaDualLow{} to \numKappaDualHigh); the LLM suggestion covered \numLLMCovered{} records with $\kappa$ = \numKappaLLM{} against the human decision, reported separately. The included records mapped to \numwordDatasets{} datasets, of which \numwordDatasetsDerived{} are derived from another public dataset and \numwordDatasetsFigureDerived{} are figure-derived collections.

\begin{figure}[H]
\centering
@@FIGURE:prisma_flow@@
\caption{PRISMA two-arm flow of records and datasets. Totals only; per-source rows are in \tabref{sources}.}\label{fig:flow}
\end{figure}

\subsection{Characteristics of sources of evidence}\label{sec:characteristics}
\tabref{datasets_main} lists every included dataset with its year, country, modality, region, subject count, annotation types, the origin of its level names, its variant labelling, licence tier, access route, quality total and fitness; the full charting records are Supplement~S1. CT is present in \numwordDatasetsCT{} datasets, plain radiographs in \numwordDatasetsXR, fluoroscopy in \numwordDatasetsFluoro{} and DXA vertebral fracture assessment in \numwordDatasetsDXA; \numwordDatasetsWholeSpine{} cover the whole spine in every image and \numwordDatasetsThoracic{} include the thoracic spine. Reported subjects summed across datasets, one per parent family, came to \numSubjectsTotal{} (persons may recur). Counts and totals by stratum are in \tabref{counts} and by year in \figref{years}.

@@TABLE:datasets_main@@

\begin{figure}[H]
\centering
@@FIGURE:datasets_by_year@@
\caption{Datasets per year of first description, stacked by modality, and cumulative reported subjects by modality.}\label{fig:years}
\end{figure}

@@TABLE:counts@@

\subsection{Critical appraisal within sources of evidence}\label{sec:appraisal_results}
Per-item scores and totals are in \tabref{quality}. The median total was \numQualityMedian{} (IQR \numQualityQOne{} to \numQualityQThree). Fit for level identification were \numwordFitLevelId{} datasets; \numwordUsableTraining{} were usable for training, \numwordFitVariant{} also labelled a variant, and \numwordFitStitching{} covered the whole spine or the thoracic spine with expert level names; the failing rule of the rest is recorded per dataset. The total was related to year with $\rho$ = \numRhoQualityYear{} and differed across access routes with $p$ = \numPKruskalRoute{} on the annotated stratum (\figref{quality}). Rubric reliability on the blind re-charted sample of \numReliabilitySample{} datasets was ICC(2,1) = \numICCTotal{} (\numICCTotalLow{} to \numICCTotalHigh) on the total; the per-item weighted $\kappa$ values are in \tabref{quality}'s note.

@@TABLE:quality@@

\begin{figure}[H]
\centering
@@FIGURE:quality@@
\caption{Quality total: histogram by access route (a) and scatter against year with a lowess line (b).}\label{fig:quality}
\end{figure}

\subsection{Results of individual sources of evidence}\label{sec:individual}
The annotation origin, variant labelling and access route of every dataset are the corresponding columns of \tabref{datasets_main}; \tabref{pooled_variants} gives every reported variant count with its exact interval, and \tabref{agreement} every reported inter-rater statistic with its comparator and category count. Per-vertebra level names were released by \numwordLevelNamesYes{} datasets; in \numwordCountingInferred{} the names were counting-inferred, which is the origin this review exists to expose, because a counting-inferred level name cannot test a counting-free method. A variant was labelled per subject with a computable prevalence in \numwordVariantExplicit{} datasets.

\subsection{Synthesis of results}\label{sec:synthesis_results}
\paragraph{Availability.} Openly downloadable were \numwordOpenDownload{} datasets; \numwordRegistration{} required registration, \numwordDUA{} a data-use agreement, \numwordRequestToAuthors{} a request to the authors, \numwordUnavailable{} were unavailable on both probe dates and \numwordCommercial{} were commercial (\tabref{availability}, \figref{availability}). Of \numwordDescribedPublic{} datasets described as public in their own text, \numwordDescribedPublicDownloads{} (\numPctDescribedPublicDownloads) had a final route of open-download. The verifier overrode the agreed machine classification in \numwordVerifierOverrides{} datasets and \numwordProbeChanged{} changed classification or status class between the two probe dates.

@@TABLE:availability@@

\begin{figure}[H]
\centering
@@FIGURE:availability@@
\caption{Access route by licence tier (mosaic) and the proportion unavailable by year band with Wilson intervals.}\label{fig:availability}
\end{figure}

\paragraph{Variant prevalence.} Eligible for pooling were \numwordPoolableLSix{} datasets for L6, \numwordPoolableTThirteen{} for T13 and \numwordPoolableLSTV{} for LSTV (\tabref{pooled_variants}, \figref{variants}). Where pooling was done, the illustrative pooled LSTV prevalence among unselected labelled datasets, by the estimator the analysis file names (\numPooledEstimator), was \numPooledLSTV{} (\numPooledLSTVLow{} to \numPooledLSTVHigh) with a prediction interval of \numPiLSTVLow{} to \numPiLSTVHigh; the prediction interval, not the point, is the finding, and enriched datasets are shown as their own stratum.

@@TABLE:pooled_variants@@

\begin{figure}[H]
\centering
@@FIGURE:variants@@
\caption{Per-dataset prevalence of L6, T13 and LSTV with exact intervals, squares scaled by subjects; illustrative pooled diamond and dashed prediction interval when at least five datasets were eligible.}\label{fig:variants}
\end{figure}

\paragraph{Inter-rater agreement.} A comparable inter-rater statistic was reported by \numwordAgreementDatasets{} datasets (\tabref{agreement}); pooling was attempted only within $\kappa$ groups meeting the plan's conditions, and the table states where it was not done and why.

@@TABLE:agreement@@

\paragraph{Coverage of the authors' grid (author-specific).} Already on the group's compute grid were \numwordOnGrid{} included datasets; \numwordReachableNotOnGrid{} reachable datasets fit for level identification were not, and they are the recipe's first-run order (\tabref{grid}).

@@TABLE:grid@@

\section{Discussion}

\subsection{Summary of evidence}\label{sec:summary}
\pending{Summary of evidence: written from the generated tables once the analysis run is frozen; each claim cites a macro}. The map answers the five questions per dataset and in aggregate: what exists and how large it is; how it is annotated and whether the level names were read or counted; whether L6, T13 and LSTV are labelled; the licence; and the route by which a reader can obtain it today. For the intended use --- a level-identification method that must be tested on expert-named, variant-labelled images --- the fit-for-level-identification list and the stitching subset are the deliverables, and the recipe reproduces their download.

\subsection{Limitations}\label{sec:limitations}
Written before the results and reported whatever they are. (1) The pooled summaries are not part of the JBI scoping design; they are secondary and illustrative, and no conclusion rests on a pooled point. (2) No access request was sent, so request-to-authors and unavailable are unrequested outcomes measured on two dates; a dataset may be obtainable by asking. (3) One reviewer screened the full title/abstract set; agreement is estimated on a stratified sample and on the whole description-review stage, so $\kappa$ describes the sample. (4) The prescreen has a measured false-negative rate per audit stratum; records it dropped that the audit did not sample are missed. (5) Several sources were reached through undocumented or HTML endpoints with named fallbacks, and their coverage is unknown; the source-overlap diagnostic shows which datasets a single source found. (6) An LLM assisted screening and drafting; its suggestions and their agreement with the human decision are published, and it made no final decision. (7) The authors maintain datasets under review; those were charted, verified and scored by another author. (8) Grid coverage is author-specific and not reproducible by a second team. (9) Raw response caches for sources whose terms restrict redistribution are published as identifiers and hashes only.

\subsection{Conclusions}\label{sec:conclusions}
\pending{Conclusions: written from the frozen results}. Every count in this report is regenerated from the recorded queries, the screening decisions, the charting records and the probes by one command, so the enumeration can be rerun on any later date and the difference reported as an update.

\section*{Funding}\label{sec:funding}
\pending{Funding statement and the funders' role}.

\section*{Competing interests}\label{sec:competing}
The authors maintain CTSpinoPelvic1K and the fluoroscopy collection named in the Introduction; both were charted, verified and scored by an author who did not create them, and the fluoroscopy collection enters only if a public describing record existed on the run date. \pending{Other competing interests}.

\section*{Data and code availability}\label{sec:availability}
The protocol and its companion files, the code (\texttt{sdr/}), the run directory with every request log, cached response (or its hash where redistribution is restricted) and \texttt{queries.json}, the screening decisions, the charting records with their quotes, the probe results, the analysis JSON files that carry every number with the files it was derived from, the generated tables and figures, and the download recipe are in the \texttt{spine\_data\_review/} directory of the repository at commit \texttt{@@GIT_COMMIT@@}; an archived copy of the frozen run is deposited with a DOI on publication. The review downloads no images and none are redistributed. Macros still pending at this render: @@PENDING_LIST@@.

\bibliographystyle{unsrtnat}
\bibliography{refs}

\appendix
\section{Verbatim search strings, endpoints, dates and counts}\label{app:search}
Generated from \texttt{queries.json} of run \runDate{}; the pilot search rows are marked \emph{pilot}, and the citation-chasing requests are collapsed to one row.
@@TABLE:search_strings@@

\section{Deviations from protocol}\label{app:deviations}
@@TABLE:deviations@@

\section{PRISMA-ScR checklist}\label{app:checklist}
@@TABLE:checklist@@

\end{document}
